\clearpage
\VerdierPhysicalPageAnchor{0212}%
\begingroup\small\setlength{\parskip}{0pt}\setlength{\emergencystretch}{1.5em}
\setlength{\abovedisplayskip}{4pt plus 2pt minus 2pt}%
\setlength{\belowdisplayskip}{4pt plus 2pt minus 2pt}%

\VerdierSemanticAnchor{ch3d-3-2-2-c}
\noindent\textit{c) 当 $n=1$ 时，$\delta^1$ 是由典范 $\delta$-函子
$\operatorname{comp}(\mathcal A)\longrightarrow D(\mathcal A)$
所定义的态射 (1.3.2)。}

\VerdierSemanticAnchor{ch3d-3-2-2-proof}
本命题由命题 (3.1.3) 的公式 (7) 得到。我们把练习留给读者，读者可参阅 [3]。

\VerdierSemanticAnchor{ch3d-3-2-3}
\medskip\noindent\textbf{3.2.3.}\quad
设
$S^n(X,Y)=0\longrightarrow Y\xrightarrow{u_n}Z_{n-1}
\xrightarrow{u_{n-1}}Z_{n-2}\longrightarrow\cdots\longrightarrow Z_0
\xrightarrow{u_0}X\longrightarrow0$
是 $E_{\mathcal A}^{n}(X,Y)$ 的一个元素。对每个满足
$0\leq p\leq n$ 的整数 $p$，给定 $\{-1,+1\}$ 的一个元素
$\epsilon_p$。以 $S_{\epsilon}^{n}(X,Y)$ 表示正合列

\VerdierSemanticAnchor{ch3d-3-2-3-1}
\begin{equation}
\resizebox{0.94\textwidth}{!}{$\displaystyle
0\longrightarrow Y\xrightarrow{\epsilon_nu_n}Z_{n-1}
\xrightarrow{\epsilon_{n-1}u_{n-1}}Z_{n-2}\longrightarrow\cdots
\longrightarrow Z_0\xrightarrow{\epsilon_0u_0}X\longrightarrow0,
$}
\tag{3.2.3.1}
\end{equation}
它是 $E_{\mathcal A}^{n}(X,Y)$ 的一个元素。于是有等式

\VerdierSemanticAnchor{ch3d-3-2-3-2}
\begin{equation}
\delta^n(S_{\epsilon}^{n}(X,Y))=
\left(\prod_{i=0}^{n}\epsilon_i\right)\delta^n(S^n(X,Y)).
\tag{3.2.3.2}
\end{equation}

\VerdierSemanticAnchor{ch3d-3-2-3-proof}
利用等价关系 (3.2.2), b) 与 $\delta^n$ 的定义 (3.2.1.1)，立即得到公式
(3.2.3.2)。

\VerdierSemanticAnchor{ch3d-3-2-4}
\medskip\noindent\textbf{3.2.4.}\quad
现在设有两个正合列：
\[
S^n(X,X')=0\longrightarrow X'\xrightarrow{u_n}Z_{n-1}
\xrightarrow{u_{n-1}}\cdots\xrightarrow{u_1}Z_0
\xrightarrow{u_0}X\longrightarrow0
\]
\[
S^m(X',X'')=0\longrightarrow X''\xrightarrow{u'_m}Z'_{m-1}
\xrightarrow{u'_{m-1}}\cdots\xrightarrow{u'_1}Z'_0
\xrightarrow{u'_0}X'\longrightarrow0.
\]
以 $S^m(X',X'')*S^n(X,X')$ 表示正合列

\VerdierSemanticAnchor{ch3d-3-2-4-1}
\begin{equation}
\begin{aligned}
S^m(X',X'')*S^n(X,X')={}&
0\longrightarrow X''\xrightarrow{u'_m}Z'_{m-1}
\xrightarrow{u'_{m-1}}\cdots\longrightarrow Z'_0
\xrightarrow{u_nu'_0}Z_{n-1}\\
&\xrightarrow{u_{n-1}}\cdots\longrightarrow Z_0
\xrightarrow{u_0}X\longrightarrow0.
\end{aligned}
\tag{3.2.4.1}
\end{equation}

\VerdierSemanticAnchor{ch3d-3-2-5}
\medskip\noindent\textbf{命题 3.2.5.}\quad
\textit{有等式}
\[
\delta^{m+n}\bigl(S^m(X',X'')*S^n(X,X')\bigr)=
(-1)^{mn}\bigl(\delta^m(S^m(X',X''))\bigr)[n]
\delta^n(S^n(X,X')).
\]

\VerdierSemanticAnchor{ch3d-3-2-5-proof-start}
来确定态射
$\bigl(\delta^m(S^m(X',X''))\bigr)[n]\delta^n(S^n(X,X'))$。
采用 (3.2.1) 的记号，有两个图式：

\VerdierSemanticAnchor{ch3d-3-2-5-first-diagram}
\[
\begin{tikzcd}[column sep=4.5em]
X' & {Z'}^\bullet
  \arrow[l,"{u(S^m(X',X''))}"']
  \arrow[r,"{v(S^m(X',X''))}"]
& X''[m]
\end{tikzcd}
\quad,
\]
\endgroup
